\documentclass[../../main.tex]{subfiles}
\graphicspath{{../../image/}{../../figure/}}
\begin{document}

\section{初边值问题}\label{section:ch06-初边值问题}

设 $\mathscr{D}$ 是 $\mathbb{R}^{m}$ 中开有界集，光滑边界 $\partial\mathscr{D}$，外单位法向场 $v$。此处考虑初边值问题
\begin{align}
\partial_{t}u(x,t)+\operatorname{div}G(u(x,t))=0,\qquad (x,t)\in\mathscr{X},
\label{align:ch06-6-9-1}
\end{align}
\begin{align}
u(x,t)=0,\qquad (x,t)\in\mathscr{B},
\label{align:ch06-6-9-2}
\end{align}
\begin{align}
u(x,0)=u_{0}(x),\qquad x\in\mathscr{D},
\label{align:ch06-6-9-3}
\end{align}
定义域为 $\mathscr{X}=\mathscr{D}\times(0,\infty)$，侧向边界 $\mathscr{B}=\partial\mathscr{D}\times(0,\infty)$。

边界条件\eqref{align:ch06-6-9-2}将在黏性消失法的框架下解释，如\cref{section:ch04-初边值问题}所释。

不等式\eqref{align:ch04-常压}启发了如下可容许弱解概念：


\begin{definition}\label{definition:ch06-6-9-1}
$\mathscr{X}$ 上的有界可测函数 $u$ 称为\eqref{align:ch06-6-9-1}、\eqref{align:ch06-6-9-2}、\eqref{align:ch06-6-9-3}以 $u_{0}\in L^{\infty}(\mathscr{D})$ 为初值的\emph{可容许弱解}，如果对每个凸熵 $\eta$（相应熵流 $Q$ 由\eqref{align:ch06-6-2-1}确定）以及 $\mathbb{R}^{m}\times[0,\infty)$ 上所有具紧支集的非负 Lipschitz 连续试验函数 $\psi$，不等式
\begin{align}
&\int_{0}^{\infty}\int_{\mathscr{D}}\Bigl{[}\partial_{t}\psi\,\eta(u)+\sum_{\alpha=1}^{m}\partial_{\alpha}\psi\,Q_{\alpha}(u)\Bigr{]}\,\mathrm{d}x\,\mathrm{d}t+\int_{\mathscr{D}}\psi(x,0)\eta(u_{0}(x))\,\mathrm{d}x\notag\\
&\qquad\geqslant\int_{0}^{\infty}\int_{\partial\mathscr{D}}\psi\,\Bigl{\{}Q_{\mathscr{B}}^{0}-\eta^{\prime}(0)[G_{\mathscr{B}}^{0}-G_{\mathscr{B}}]\Bigr{\}}\,\mathrm{d}\mathscr{H}^{m-1}(x)\,\mathrm{d}t
\label{align:ch06-6-9-4}
\end{align}
成立。$G_{\mathscr{B}}$ 表示 $G$ 的法向分量在 $\mathscr{B}$ 上的迹，$G_{\mathscr{B}}^{0}$ 与 $Q_{\mathscr{B}}^{0}$ 分别表示 $G(0)v$ 与 $Q(0)v$。
\end{definition}

注意\eqref{align:ch06-6-9-4}蕴含 $\partial_{t}\eta+\operatorname{div}Q\leqslant 0$，特别地 $\partial_{t}u+\operatorname{div}G=0$，故 $Q$ 与 $G$ 的法向分量在 $\mathscr{B}$ 上的迹 $Q_{\mathscr{B}}$、$G_{\mathscr{B}}$ 良定义。此外，\eqref{align:ch04-PBG}在 $\mathscr{B}$ 上以形式

\begin{align}
Q_{\mathscr{B}}-Q_{\mathscr{B}}^{0}-\eta^{\prime}(0)[G_{\mathscr{B}}-G_{\mathscr{B}}^{0}]\geqslant 0
\label{align:ch06-6-9-5}
\end{align}
成立。

以技术上的复杂化为代价，但无本质困难，特殊边界条件 $u=0$ 可替换为 $u=\hat{u}(x,t)$，其中 $\hat{u}$ 为任一充分光滑函数。

\cref{definition:ch06-6-9-1}的合理性由下述定理提供。

\begin{theorem}\label{theorem:ch06-6-9-2}
对每个 $u_{0}\in L^{\infty}(\mathscr{D})$，存在\eqref{align:ch06-6-9-1}、\eqref{align:ch06-6-9-2}、\eqref{align:ch06-6-9-3}的唯一可容许弱解 $u$，且
\begin{align}
u(\cdot,t)\in C^{0}([0,\infty);L^{1}(\mathscr{D})).
\label{align:ch06-6-9-6}
\end{align}
此外，若 $u_{0}\in BV(\mathscr{D})$，则 $u\in BV_{\mathrm{loc}}(\mathscr{X})$。
\end{theorem}

在通过证明上述定理来建立解的存在性之前，先借\cref{theorem:ch06-6-2-3}的下述类比式证明唯一性与稳定性。

\begin{theorem}\label{theorem:ch06-6-9-3}
设 $u$ 与 $\bar{u}$ 是\eqref{align:ch06-6-9-1}、\eqref{align:ch06-6-9-2}以 $u_{0}$、$\bar{u}_{0}$ 为相应初值的可容许弱解。则对任意 $t>0$，
\begin{align}
\int_{\mathscr{D}}[u(x,t)-\bar{u}(x,t)]^{+}\,\mathrm{d}x\leqslant\int_{\mathscr{D}}[u_{0}(x)-\bar{u}_{0}(x)]^{+}\,\mathrm{d}x,
\label{align:ch06-6-9-7}
\end{align}
\begin{align}
\|u(\cdot,t)-\bar{u}(\cdot,t)\|_{L^{1}(\mathscr{D})}\leqslant\|u_{0}(\cdot)-\bar{u}_{0}(\cdot)\|_{L^{1}(\mathscr{D})}.
\label{align:ch06-6-9-8}
\end{align}
此外，若
\begin{align}
u_{0}(x)\leqslant\bar{u}_{0}(x),\qquad\text{a.e.\ 在 }\mathscr{D}\text{ 上},
\label{align:ch06-6-9-9}
\end{align}
则
\begin{align}
u(x,t)\leqslant\bar{u}(x,t),\qquad\text{a.e.\ 在 }\mathscr{X}\text{ 上}.
\label{align:ch06-6-9-10}
\end{align}
\end{theorem}

\begin{proof}
在简化假设下证明梗概：设 $u$ 与 $\bar{u}$ 在 $\mathscr{B}$ 上均取得强迹 $u_{\mathscr{B}}$、$\bar{u}_{\mathscr{B}}$，此时 $G$ 与 $Q$ 的法向分量在 $\mathscr{B}$ 上的迹由普通复合得到：
\begin{align}
G_{\mathscr{B}}=G(u_{\mathscr{B}})v,\ Q_{\mathscr{B}}=Q(u_{\mathscr{B}})v,\ \bar{G}_{\mathscr{B}}=G(\bar{u}_{\mathscr{B}})v,\ \bar{Q}_{\mathscr{B}}=Q(\bar{u}_{\mathscr{B}})v.
\label{align:ch06-6-9-11}
\end{align}
当 $u$、$\bar{u}$ 为 $BV$ 函数，或当 $u$、$\bar{u}$ 仅属于 L$^{\infty}$ 而 $G$ 线性非退化时，上述假设成立；见\cref{theorem:ch06-6-8-3}。

重走\cref{theorem:ch06-6-2-3}证明的步骤，采用由\eqref{align:ch06-6-2-5}定义的同一熵--熵流对 $(\eta(u;\bar{u}),Q(u;\bar{u}))$ 与由\eqref{align:ch06-6-2-16}给出的同一试验函数 $\phi(x,t,\bar{x},\bar{t})$。但现在在 $\mathscr{D}\times[0,\infty)$ 而非 $\mathbb{R}^{m}\times[0,\infty)$ 上积分，并用\eqref{align:ch06-6-9-4}代替\eqref{align:ch06-6-2-3}。于是代替\eqref{align:ch06-6-2-21}得到
\begin{align}
&\int_{0}^{\infty}\int_{\mathscr{D}}\Bigl{\{}\partial_{t}\psi\,\eta(u;\bar{u})+\sum_{\alpha=1}^{m}\partial_{\alpha}\psi\,Q_{\alpha}(u;\bar{u})\Bigr{\}}\,\mathrm{d}x\,\mathrm{d}t+\int_{\mathscr{D}}\psi(x,0)\eta(u_{0}(x);\bar{u}_{0}(x))\,\mathrm{d}x\notag\\
&\qquad\geqslant\int_{0}^{\infty}\int_{\partial\mathscr{D}}\psi\,\operatorname{sgn}[u_{\mathscr{B}}-\bar{u}_{\mathscr{B}}]^{+}[G_{\mathscr{B}}-\bar{G}_{\mathscr{B}}]\,\mathrm{d}\mathscr{H}^{m-1}(x)\,\mathrm{d}t.
\label{align:ch06-6-9-12}
\end{align}
由边界条件\eqref{align:ch06-6-9-5}可知，\eqref{align:ch06-6-9-12}右端的积分为非负。事实上，被积函数在 $u_{\mathscr{B}}\leqslant\bar{u}_{\mathscr{B}}$ 处为零，在 $u_{\mathscr{B}}>\bar{u}_{\mathscr{B}}$ 处与 $G_{\mathscr{B}}-\bar{G}_{\mathscr{B}}$ 同号。在后一情形，分别考察以下三个子情形：
\begin{enumerate}
\item[(a)] $u_{\mathscr{B}}>\bar{u}_{\mathscr{B}}\geqslant 0$：对解 $u$ 与熵--熵流对 $(\eta(u;\bar{u}_{\mathscr{B}}),Q(u;\bar{u}_{\mathscr{B}}))$ 写出\eqref{align:ch06-6-9-5}，得 $G_{\mathscr{B}}\geqslant\bar{G}_{\mathscr{B}}$。
\item[(b)] $0\geqslant u_{\mathscr{B}}>\bar{u}_{\mathscr{B}}$：对解 $\bar{u}$ 与熵--熵流对 $(\eta(u_{\mathscr{B}};\bar{u}),Q(u_{\mathscr{B}};\bar{u}))$ 写出\eqref{align:ch06-6-9-5}，同样得 $G_{\mathscr{B}}\geqslant\bar{G}_{\mathscr{B}}$。
\item[(c)] $u_{\mathscr{B}}>0>\bar{u}_{\mathscr{B}}$：对解 $u$ 与熵--熵流对 $(\eta(u;0),Q(u;0))$ 写出\eqref{align:ch06-6-9-5}，得 $G_{\mathscr{B}}\geqslant G_{\mathscr{B}}^{0}$。类似地，对解 $\bar{u}$ 与熵--熵流对 $(\eta(0;\bar{u}),Q(0;\bar{u}))$ 写出\eqref{align:ch06-6-9-5}，得 $\bar{G}_{\mathscr{B}}\leqslant G_{\mathscr{B}}^{0}$。特别地，$G_{\mathscr{B}}\geqslant\bar{G}_{\mathscr{B}}$。
\end{enumerate}

对试验函数 $\psi(x,\tau)=\chi(x)\omega(\tau)$ 应用\eqref{align:ch06-6-9-12}，其中 $x\in\mathscr{D}$ 时 $\chi(x)=1$，$\omega$ 由\eqref{align:ch05-cone-convex}定义。

由于\eqref{align:ch06-6-9-12}右端非负，得
\begin{align}
\frac{1}{\varepsilon}\int_{t}^{t+\varepsilon}\int_{\mathscr{D}}[u(x,\tau)-\bar{u}(x,\tau)]^{+}\,\mathrm{d}x\,\mathrm{d}\tau\leqslant\int_{\mathscr{D}}[u_{0}(x)-\bar{u}_{0}(x)]^{+}\,\mathrm{d}x.
\label{align:ch06-6-9-13}
\end{align}
令 $\varepsilon\downarrow 0$，即得\eqref{align:ch06-6-9-7}。进而\eqref{align:ch06-6-9-7}直接蕴含定理其余论断。证明完成。
\end{proof}

下一步用黏性消失法构造\eqref{align:ch06-6-9-1}、\eqref{align:ch06-6-9-2}、\eqref{align:ch06-6-9-3}的解。考虑抛物方程组
\begin{align}
\partial_{t}u(x,t)+\operatorname{div}G(u(x,t))=\mu\Delta u(x,t),\qquad (x,t)\in\mathscr{X},
\label{align:ch06-6-9-14}
\end{align}
带边界条件\eqref{align:ch06-6-9-2}与初始条件\eqref{align:ch06-6-9-3}。对任意 $\mu>0$，\eqref{align:ch06-6-9-14}、\eqref{align:ch06-6-9-2}、\eqref{align:ch06-6-9-3}存在唯一解 $u_{\mu}$，在 $\bar{\mathscr{D}}\times(0,\infty)$ 上光滑。由最大值原理，
\begin{align}
|u_{\mu}(x,t)|\leqslant\sup|u_{0}(\cdot)|,\qquad x\in\mathscr{D},\ t\in(0,\infty).
\label{align:ch06-6-9-15}
\end{align}

重走\cref{theorem:ch06-6-3-2}证明的步骤，只是现在把\eqref{align:ch06-6-3-10}在 $\mathscr{D}\times(s,t)$ 而非 $\mathbb{R}^{m}\times(s,t)$ 上积分，易得

\begin{theorem}\label{theorem:ch06-6-9-4}
设 $u_{\mu}$ 与 $\bar{u}_{\mu}$ 是\eqref{align:ch06-6-9-14}、\eqref{align:ch06-6-9-2}以 $u_{0}$、$\bar{u}_{0}$ 为相应初值的解。则对任意 $t>0$，
\begin{align}
\int_{\mathscr{D}}[u_{\mu}(x,t)-\bar{u}_{\mu}(x,t)]^{+}\,\mathrm{d}x\leqslant\int_{\mathscr{D}}[u_{0}(x)-\bar{u}_{0}(x)]^{+}\,\mathrm{d}x,
\label{align:ch06-6-9-16}
\end{align}
\begin{align}
\|u_{\mu}(\cdot,t)-\bar{u}_{\mu}(\cdot,t)\|_{L^{1}(\mathscr{D})}\leqslant\|u_{0}(\cdot)-\bar{u}_{0}(\cdot)\|_{L^{1}(\mathscr{D})}.
\label{align:ch06-6-9-17}
\end{align}
此外，若
\begin{align}
u_{0}(x)\leqslant\bar{u}_{0}(x),\qquad\text{a.e.\ 在 }\mathscr{D}\text{ 上},
\label{align:ch06-6-9-18}
\end{align}
则
\begin{align}
u_{\mu}(x,t)\leqslant\bar{u}_{\mu}(x,t),\qquad (x,t)\in\mathscr{D}\times(0,\infty).
\label{align:ch06-6-9-19}
\end{align}
\end{theorem}

现证明\eqref{align:ch06-6-9-14}、\eqref{align:ch06-6-9-2}、\eqref{align:ch06-6-9-3}的解族 $\{u_{\mu}:\mu>0\}$ 在 $L^{1}$ 中相对紧。

\begin{lemma}\label{lemma:ch06-6-9-5}
设 $u_{\mu}$ 为\eqref{align:ch06-6-9-14}、\eqref{align:ch06-6-9-2}、\eqref{align:ch06-6-9-3}以 $u_{0}\in L^{\infty}(\mathscr{D})\cap W^{2,1}(\mathscr{D})$ 为初值的解。则对任意 $t>0$，
\begin{align}
\|\partial_{t}u_{\mu}(\cdot,t)\|_{L^{1}(\mathscr{D})}\leqslant c_{0}\|u_{0}(\cdot)\|_{W^{1,1}(\mathscr{D})}+\mu\,\|u_{0}(\cdot)\|_{W^{2,1}(\mathscr{D})},
\label{align:ch06-6-9-20}
\end{align}
\begin{align}
\sum_{\beta=1}^{m}\|\partial_{\beta}u_{\mu}(\cdot,t)\|_{L^{1}(\mathscr{D})}\leqslant a(t)\|u_{0}(\cdot)\|_{W^{1,1}(\mathscr{D})}+\mu\,b(t)\|u_{0}(\cdot)\|_{W^{2,1}(\mathscr{D})},
\label{align:ch06-6-9-21}
\end{align}
其中 $c_{0}$ 与连续函数 $a(t)$、$b(t)$ 不依赖于 $\mu$。
\end{lemma}

\begin{proof}
对 $h>0$，对两解 $u_{\mu}(x,t)$（初值 $u_{0}(x)$）与 $\bar{u}_{\mu}(x,t)=u_{\mu}(x,t+h)$（初值 $\bar{u}_{0}(x)=u_{\mu}(x,h)$）应用\eqref{align:ch06-6-9-17}。除以 $h$ 再令 $h\downarrow 0$，得 $\|\partial_{t}u_{\mu}(\cdot,t)\|_{L^{1}(\mathscr{D})}\leqslant\|\partial_{t}u_{\mu}(\cdot,0)\|_{L^{1}(\mathscr{D})}$，借助\eqref{align:ch06-6-9-14}即得\eqref{align:ch06-6-9-20}。

不能用同样程序估计空间导数，因为空间方向平移不再把解变为解。故须用不同论证。

对 $\varepsilon>0$，定义函数
\begin{align}
\eta_{\varepsilon}(w)=
\begin{cases}
-w-\varepsilon, & -\infty<w\leqslant-2\varepsilon,\\
\dfrac{w^{2}}{4\varepsilon}, & -2\varepsilon<w\leqslant 2\varepsilon,\\[6pt]
w-\varepsilon, & 2\varepsilon<w<\infty.
\end{cases}
\label{align:ch06-6-9-22}
\end{align}
令 $w=\partial_{\beta}u_{\mu}$，关于 $x_{\beta}$ 微分\eqref{align:ch06-6-9-14}，以 $\eta_{\varepsilon}^{\prime}(w)$ 乘所得方程并在 $\mathscr{D}$ 上积分。分部积分后得
\begin{align}
\frac{\mathrm{d}}{\mathrm{d}t}\int_{\mathscr{D}}\eta_{\varepsilon}(w)\,\mathrm{d}x
&=\int_{\mathscr{D}}[\eta_{\varepsilon}(w)-\eta_{\varepsilon}^{\prime}(w)w]\operatorname{div}G^{\prime}(u_{\mu})\,\mathrm{d}x\notag\\
&\quad-\mu\int_{\mathscr{D}}\eta_{\varepsilon}^{\prime\prime}(w)|\nabla w|^{2}\,\mathrm{d}x+\int_{\partial\mathscr{D}}\Bigl{[}\mu\eta_{\varepsilon}^{\prime}(w)\frac{\partial w}{\partial v}-\eta_{\varepsilon}(w)G^{\prime}(0)v\Bigr{]}\,\mathrm{d}\mathscr{H}^{m-1}(x).
\label{align:ch06-6-9-23}
\end{align}
当 $\varepsilon\downarrow 0$，\eqref{align:ch06-6-9-23}左端被积函数趋于 $|w|$。右端第一个积分为 $O(\varepsilon)$，第二个积分非负。为估计 $\partial\mathscr{D}$ 上的积分，注意由于 $u_{\mu}$ 在边界上为零，$\partial_{\alpha}u_{\mu}=\frac{\partial u_{\mu}}{\partial v}v_{\alpha}$（$\alpha=1,\cdots,m$）。特别地，$w=\frac{\partial u_{\mu}}{\partial v}v_{\beta}$。于是\eqref{align:ch06-6-9-14}蕴含 $\frac{\partial u_{\mu}}{\partial v}G^{\prime}(0)v=\mu\Delta u_{\mu}$。最后显然有 $\frac{\partial w}{\partial v}=\frac{\partial^{2}u_{\mu}}{\partial v^{2}}v_{\beta}+O(1)\frac{\partial u_{\mu}}{\partial v}$，且 $\Delta u_{\mu}=\frac{\partial^{2}u_{\mu}}{\partial v^{2}}+O(1)\frac{\partial u_{\mu}}{\partial v}$。于是
\begin{align}
\mu\eta_{\varepsilon}^{\prime}(w)\frac{\partial w}{\partial v}-\eta_{\varepsilon}(w)G^{\prime}(0)v
=\mu\Bigl{[}\eta_{\varepsilon}^{\prime}(w)-\frac{\eta_{\varepsilon}(w)}{w}\Bigr{]}\frac{\partial^{2}u_{\mu}}{\partial v^{2}}v_{\beta}+O(1)\mu\frac{\partial u_{\mu}}{\partial v},
\label{align:ch06-6-9-24}
\end{align}
当 $\varepsilon\downarrow 0$ 时趋于 $O(1)\mu\frac{\partial u_{\mu}}{\partial v}$。因而令 $\varepsilon\downarrow 0$，\eqref{align:ch06-6-9-23}给出
\begin{align}
\frac{\mathrm{d}}{\mathrm{d}t}\int_{\mathscr{D}}|\partial_{\beta}u_{\mu}|\,\mathrm{d}x\leqslant c\int_{\partial\mathscr{D}}\mu\Bigl{|}\frac{\partial u_{\mu}}{\partial v}\Bigr{|}\,\mathrm{d}\mathscr{H}^{m-1}(x)\leqslant c^{\prime}\int_{\mathscr{D}}\mu|\Delta u_{\mu}|\,\mathrm{d}x.
\label{align:ch06-6-9-25}
\end{align}
对 $\beta=1,\cdots,m$ 求和\eqref{align:ch06-6-9-25}，并用 $\partial_{t}u_{\mu}+\operatorname{div}G(u_{\mu})$ 代替 $\mu\Delta u_{\mu}$。利用\eqref{align:ch06-6-9-15}、\eqref{align:ch06-6-9-20}并应用 Gronwall 不等式，即得\eqref{align:ch06-6-9-21}。证明完成。
\end{proof}

\begin{proof}[\cref{theorem:ch06-6-9-2}的证明]
先设 $u_{0}\in L^{\infty}(\mathscr{D})\cap W^{2,1}(\mathscr{D})$。由\cref{lemma:ch06-6-9-5}，\eqref{align:ch06-6-9-14}、\eqref{align:ch06-6-9-2}、\eqref{align:ch06-6-9-3}的解族 $\{u_{\mu}:\mu>0\}$ 在 $L^{1}(\mathscr{D}\times(0,T))$ 中对任一 $T>0$ 相对紧。于是回顾\eqref{align:ch06-6-9-15}，可提取序列 $\{u_{\mu_{k}}\}$（$\mu_{k}\downarrow 0$，$k\to\infty$），它在 $\mathscr{D}\times(0,\infty)$ 上一致有界且几乎处处收敛于某函数 $u$。正如\cref{section:ch04-初边值问题}所示，$u$ 满足\eqref{align:ch06-6-9-4}，因而是\eqref{align:ch06-6-9-1}、\eqref{align:ch06-6-9-2}、\eqref{align:ch06-6-9-3}的唯一解。

特别地，整个族 $\{u_{\mu}:\mu>0\}$ 当 $\mu\downarrow 0$ 时收敛于 $u$。此外，由\eqref{align:ch06-6-9-20}、\eqref{align:ch06-6-9-21}，$u\in BV_{\mathrm{loc}}(\mathscr{D}\times(0,\infty))$，且对任一 $T>0$，
\begin{align}
TV_{\mathscr{D}\times(0,T)}u\leqslant c(T)\|u_{0}\|_{W^{1,1}(\mathscr{D})}.
\label{align:ch06-6-9-26}
\end{align}
此外，$u$ 由\eqref{align:ch06-6-9-15}继承最大值原理：$|u(x,t)|\leqslant\sup|u_{0}(\cdot)|$。

现设 $u_{0}\in L^{\infty}(\mathscr{D})$。构造 $L^{\infty}(\mathscr{D})\cap W^{2,1}(\mathscr{D})$ 中函数列 $\{u_{0n}\}$，使 $\|u_{0n}\|_{L^{\infty}(\mathscr{D})}\leqslant\|u_{0}\|_{L^{\infty}(\mathscr{D})}$ 且 $u_{0n}\to u_{0}$（在 $L^{1}(\mathscr{D})$ 中）。由\eqref{align:ch06-6-9-8}，以 $u_{0n}$ 为初值的\eqref{align:ch06-6-9-1}、\eqref{align:ch06-6-9-2}的可容许解序列 $\{u_{n}\}$ 在 $L^{1}$ 中收敛于某函数 $u$，后者即\eqref{align:ch06-6-9-1}、\eqref{align:ch06-6-9-2}、\eqref{align:ch06-6-9-3}的可容许解。此外，当 $u_{0}\in BV(\mathscr{D})$ 时，序列 $\{u_{0n}\}$ 可附加构造使 $\|u_{0n}\|_{W^{1,1}(\mathscr{D})}\leqslant C[TV_{\mathscr{D}}u_{0}+\|u_{0}\|_{L^{\infty}(\mathscr{D})}]$；此时\eqref{align:ch06-6-9-26}蕴含对任一 $T>0$，$u\in BV(\mathscr{D}\times(0,T))$。证明完成。
\end{proof}

\end{document}
